\textbf{1.\ \textsc{Moving lens} (6 points)} --- \textit{Eppu Leinonen.}
The centre of a thin converging lens moves along a circle while the
orientation of the lens remains fixed; the optical axis of the lens lies in
the plane of the circle (the plane of the figure). Two fixed points $P$ and
$Q$, also in this plane, are imaged by the lens; the images are always real.
As the lens moves, each image traces a closed curve in the plane of the
figure: $P \to P'$, $Q \to Q'$. The two points and both curves are shown in
the figure.

\begin{center}\includegraphics[width=0.5\linewidth]{figures/NBPhO_2026_Q1_fig1.png}\end{center}

\textbf{i)} \textit{(2 points)} Construct the circle along which the lens
centre moves.

\textbf{ii)} \textit{(4 points)} Determine the direction of the lens plane,
i.\,e.\ construct a line parallel to the lens.

In the construction tasks, draw the construction on the provided sheet with
the figure, detail the steps of your construction, and explain why your
construction works.
